\documentclass{beamer}

\mode<presentation>
{
  \usetheme{Singapore}      % or try Darmstadt, Madrid, Warsaw, ...
  \usecolortheme{crane} % or try albatross, beaver, crane, ...
  \usefonttheme{serif}  % or try serif, structurebold, ...
  \setbeamertemplate{navigation symbols}{}
  \setbeamertemplate{caption}[numbered]
} 

\usepackage[autostyle=true]{csquotes}
\usepackage{xeCJK}
\usepackage[T1]{fontenc}
\usepackage{xcolor}
\usepackage{setspace}
\setmainfont[AutoFakeSlant=0.2]{Charis SIL}
\setCJKmainfont[AutoFakeSlant=0.2]{Noto Serif CJK HK}
\setCJKmonofont[AutoFakeSlant=0.2]{Noto Serif CJK HK}
\usepackage{polyglossia}
\newfontfamily\thaifont[Script=Thai]{Noto Serif Thai}
\setmainlanguage{english}
\setotherlanguages{thai}
\usepackage{csquotes}
\usepackage{array}
\usepackage{makecell}
\usepackage{changepage}
\usepackage{textcomp}
\usepackage{multirow}

\title[Pearl River Sinitic]{Yuè \& the other Sinitic Languages of the Pearl River Region --- Unit 4.1}
\subtitle{13ᵗʰ EACL Summer School, Universität Stuttgart}
\author{Hilário de Sousa 蘇沙}
\institute{EHESS–CRLAO \\ {\tiny membre associé}}
\date{Wednesday 30 July 2026 \\ 丙午六月十七 \\ {\footnotesize \texttt{hilario@bambooradical.com}}}

\begin{document}

\begin{frame}
  \titlepage
\end{frame}

\section{N–S \& Viet}

\begin{frame}[fragile]{North–South and Vietnam}

\begin{columns}
\begin{column}{0.8\textwidth}
\includegraphics[width=1\linewidth]{watershed1}
\end{column}
\begin{column}{0.25\textwidth}
靈渠 (秦 214 BCE near 桂林)\\
\medskip
鬼門關 (鬱林/ 玉林 \& 北流)\\
合浦\\
\medskip
梅關道 (唐 716 CE) \& 珠璣巷\\
\medskip
Pinghua: 4m; Zhuang 14m; Yue 68m\\ 
\end{column}
\end{columns}

\vfill \setstretch{0.5} {\tiny de Sousa, Hilário. 2020. On Pinghua, and Yue: Some historical and linguistic perspectives. Crossroads: An Interdisciplinary Journal of Asian Interactions 19(2). 257–295. \url{https://doi.org/10.1163/26662523-bja10004}}\par

\end{frame}

\begin{frame}[fragile]{North–South and Vietnam}

Chinese loanwords in Korean \& Japanese:\\
Approximation of the prestigious variety of Chinese at the time\\
Transmitted through scholars who went to China\\
\bigskip
Chinese loanwords in Vietnamese:\\
Not approximation of a prestigious variety of Chinese\\
Most-likely scenario:\\
Local Sinitic variety/ies spoken by the Chinese community in the Red River Delta region, i.e. \enquote{Annamese Middle Chinese} (Phan 2025; Phan \& de Sousa 2022), Fujianese (Whitmore 2020)\\

\vfill \setstretch{0.5} {\tiny Phan, John D. 2025. \textit{Lost tongues of the Red River: Annamese Middle Chinese \& the origins of the Vietnamese language}. Cambridge MA: Harvard University Press.\\
Phan, John D. \& de Sousa, Hilário. 2022. Southwestern Middle Chinese: Preliminary evidence from Hunan, Guangxi, and Sino-Vietnamese. In Phan, Trang \& Phan, John D. \& Alves, Mark J (eds.), \textit{Vietnamese linguistics: State of the field} (Journal of the Southeast Asian Linguistics Society Special Publication 9), 59–79. Honolulu: University of Hawai‘i Press.\\
Whitmore, John K. 2020. The Sông Cái (Red River) Delta, the Chinese diaspora, and the Trần/Chen clan of Đại Việt. \textit{Crossroads: An interdisciplinary journal of Asian interactions} 19. 210–232.}\par

\end{frame}

\begin{frame}[fragile]{North–South and Vietnam}

{\includegraphics[width=1.1\textwidth]{dynastiespearlred}}

\end{frame}

\begin{frame}[fragile]{North–South and Vietnam}

Alves (2018)\\
\medskip
\begin{tabular}{r|c c c c}
    Old Chin. & son. end. & -ʔ & -s & -p -t -k\\ \hline
    Middle Chin. & 平 A & 上 B & 去 C & 入 D\\
\end{tabular}

\bigskip
\begin{tabular}{r|c c c c}
    Proto-Vietic & son. end. & -ʔ & -s -h & -p -t -c -k\\ \hline
    (native) Viet. & ngang & sắc & hỏi & sắc\\
     & huyền & nặng & ngã & nặng\\ 
\end{tabular}

\bigskip
{\footnotesize
\begin{tabular}{r|c c|c c|c c|c c}
    P-Viet. & \multicolumn{2}{c|}{son. end.} & \multicolumn{2}{c|}{-ʔ} & \multicolumn{2}{c|}{-s -h} & \multicolumn{2}{c}{-p -t -c -k}\\ \hline
    P-Viet. & *ciːm & *veːr & *ʔa-cɔːʔ & *r-koːʔ & *cɛh & *muːs & *-suk & *məc\\ 
    Viet. & chim & về & chó & gạo & chẻ & mũi & tóc & mật\\ \hline
    gloss & bird & turn & dog & husked & to & nose & hair & gall\\
     & & back & & rice & split & & & \\
\end{tabular}
}

\vfill \setstretch{0.5} {\tiny Alves, Mark J. 2018. Early Sino-Vietnamese Lexical Data and the Relative Chronology of Tonogenesis in Chinese and Vietnamese. \textit{Bulletin of Chinese Linguistics} 11(1–2). 3–33. \url{https://doi.org/10.1163/2405478X-01101007}}\par

\end{frame}

\begin{frame}[fragile]{North–South and Vietnam}

Terminologies surrounding \enquote{Sino-Vietnamese}

\bigskip
{\footnotesize
\begin{tabular}{l l l}
 & common & Phan (2025)\\ \hline
0 CE (Han dyn.; LOC) & \multirow{2}{*}{Old-SV/\enquote{Non-SV}} & \multirow{2}{*}{Early SV}\\ \cline{1-1}
601 (Sui dyn.; EMC) & & \\ \hline
1000 (Viet. indep.; LMC) & Sino-Vietnamese & Late SV\\
literary pronunciation & 漢越音 & \\
\textit{Red River Chinese shifting} & âm Hán-Việt & \\
\textit{to Việt-Mường} & & \\ \hline
later, \textit{e.g. recent loans} & & Recent SV \\
\textit{from Canto./Teochew} & & \\
\textit{Sinitic neologisms} & & \\
\end{tabular}
}

\medskip
Examples of Sinitic 文白異讀\\ 
Mandarin 血 Lit. xuè e.g. 血統, Col. xiě e.g. 流血\\
Cantonese 正 Lit. zing3 e.g. 正統, Col. zeng3 e.g. 正！\\ 


\end{frame}

\begin{frame}[fragile]{North–South and Vietnam}

\begin{tabular}{l|c c c c}
MC tone & 平 A & 上 B & 去 C & 入 D\\ \hline 
\multirow{2}{*}{ESV} & \multirow{3}{*}{ngang huyền} & \multirow{2}{*}{sắc nặng} & hỏi ngã & \multirow{3}{*}{sắc nặng}\\
 & & & ngang huyền & \\ \cline{1-1} \cline{3-4}
LSV & & hỏi ngã & sắc nặng & \\
\end{tabular}

\bigskip
\begin{tabular}{r|l l}
 & ESV & LSV \\ \hline
紙 & giấy \enquote*{paper}& chỉ (e.g. chứng chỉ 證紙)\\
網 & mạng \enquote*{network} & võng \enquote*{hammock}\\
歲 & tuổi \enquote*{age} & tuế (e.g. vạn tuế 萬歲)\\
字 & chữ \enquote*{character} & tự (e.g. văn tự 文字 \enquote*{contract})\\ 
\end{tabular}

\end{frame}

\begin{frame}[fragile]{North–South and Vietnam}

\begin{tabular}{l|c c c c}
MC tone & 平 A & 上 B & 去 C & 入 D\\ \hline 
\multirow{2}{*}{ESV} & \multirow{3}{*}{ngang huyền} & \multirow{2}{*}{sắc nặng} & hỏi ngã & \multirow{3}{*}{sắc nặng}\\
 & & & ngang huyền & \\ \cline{1-1} \cline{3-4}
LSV & & hỏi ngã & sắc nặng & \\
\end{tabular}

\bigskip
\begin{tabular}{r|l l l}
 & ESV (LOC) & ESV (EMC) & LSV (LMC)\\ \hline 
帶 & dải & đai & đái (\~{}đới)\\
e.g. & dải đất \enquote*{strip of land} & đai đen \enquote*{black belt} & ôn đới 溫帶\\ \hline
共 & cũng & cùng & cộng\\
e.g. & \enquote*{also} & \enquote*{with} & tổng cộng 總共\\

\end{tabular}
\end{frame}

\begin{frame}[fragile]{North–South and Vietnam}

\begin{tabular}{l|c c c c}
MC tone & 平 A & 上 B & 去 C & 入 D\\ \hline 
\multirow{2}{*}{ESV} & \multirow{3}{*}{ngang huyền} & \multirow{2}{*}{sắc nặng} & hỏi ngã & \multirow{3}{*}{sắc nặng}\\
 & & & ngang huyền & \\ \cline{1-1} \cline{3-4}
LSV & & hỏi ngã & sắc nặng & \\
\end{tabular}

\bigskip

\begin{tabular}{l l l}
LOC: 帶 *C.tˤa[t]-s & > ESV -s (no tone) & > later: dải \textsuperscript{hỏi}\\ \hline
EMC: 帶 taj\textsuperscript{C} & > ESV -0 (no tone) & > later: đai \textsuperscript{ngang}\\ \hline
LMC: 帶 taj\textsuperscript{C} & & > LSV: đái \textsuperscript{sắc}\\ 



\end{tabular}
\end{frame}

\end{document}
